% cross_chain_gamma.tex -- round 425, the cross-chain gamma
% (reviewer DCCXCI).
% B2_LE_S_SATZ / QN_BRIDGES_BW / DEAD_NEEDS_SLACK / COFINAL_OPEN.
% Builds with pdflatex.
% Companion machine check: verify_cross_chain.py.
% Discovery: experiments/tfpt-discovery/cross_chain_gamma_probe.py.
% NO RH CLAIM.  Research documentation, not a theorem of RH.
\documentclass[11pt]{article}

\usepackage[a4paper,margin=2.7cm]{geometry}
\usepackage{amsmath,amssymb,amsthm}
\usepackage{microtype}
\usepackage{booktabs}
\usepackage{xcolor}
\usepackage[colorlinks=true,linkcolor=blue!50!black,urlcolor=blue!50!black]{hyperref}

\newtheorem{lemma}{Lemma}
\newtheorem{prop}{Proposition}
\theoremstyle{remark}
\newtheorem{remark}{Remark}

\title{\bfseries Cross-chain $\gamma$\\[2mm]
\large $\|b\|^2\le S$ is a kernel comparison;
$\|b\|^2<B_w$ waits on $q_N<1$}
\author{Stefan Hamann}
\date{August 29, 2026}

\begin{document}
\maketitle

\begin{abstract}\noindent
Round 424 identified $b$ as the $\mu$-Parseval
coefficient of the signed border $\sigma$.
Both $\|b\|^2$ and $S$ are the $N$-Christoffel
energies of the same $\sigma$, in the
$\mu$-ONB and the $\mu-\nu$-ONB of
$P_{<N}$.
On that space
$\langle p,p\rangle_\mu\ge\langle p,p\rangle_{\mu-\nu}$
as soon as $\nu$ is a positive measure and the
signed form is definite (tilde chain ok).
Reproducing kernels then satisfy
$K^{\mu-\nu}\succeq K^\mu$, hence
$\|b\|^2\le S$.
$B_w=S_{N-2}+\tfrac57$ is $S$ exactly when
$q_N=1$.
The posted $\|b\|^2<B_w$ is this SATZ plus
$q_N<1$.
No RH claim.
\end{abstract}

\medskip\noindent
\textbf{Status.}\enspace
Lemma~\ref{lem:cross} is \textbf{reduced}
(cofinal $q_N<1$ remains).
Proposition~\ref{prop:kernel} is a finite theorem.
On the $42$-rung $q_N<1$ census it
\emph{composes} to $\gamma<1$.
Ausgang \texttt{B2\_LE\_S\_SATZ / QN\_BRIDGES\_BW / DEAD\_NEEDS\_SLACK / COFINAL\_OPEN}.
No $L^{*}$ claim.  No $R^{\dagger}$ claim.  No RH claim.

\section{The lemma}\label{sec:lem}

\begin{lemma}[cross-chain $\gamma$]\label{lem:cross}
$\|b\|^2<B_w=S_{N-2}+\tfrac57$ as a
cross-chain theorem, so $\gamma<1$ and
$\mathrm{den}=1+\gamma-v_t\cdot s<2$
even at $v_t\cdot s=0$.
\end{lemma}

\noindent
\textbf{Verdict: REDUZIERT} as posted.
The sharp theorem is $\|b\|^2\le S$
(Proposition~\ref{prop:kernel}).
$S<B_w$ iff $q_N<1$.
Cofinal $q_N<1$ is the remaining gap
(already the $v960$ census, not cofinal).

\section{The common majorant}\label{sec:maj}

Let $P_{<N}$ carry the two inner products
\[
  \langle p,p\rangle_\mu=\int p^2\,d\mu,\qquad
  \langle p,p\rangle_{\mu-\nu}
  =\int p^2\,d\mu-\int p^2\,d\nu.
\]
Whenever the tilde chain exists
($h_k>0$ through the window),
the second form is positive definite on
$P_{<N}$.
Since $\nu$ is a positive measure,
$\langle p,p\rangle_\mu\ge\langle p,p\rangle_{\mu-\nu}$.
The reproducing kernels of a finite-dimensional
space then satisfy
$K_N^{\mu-\nu}\succeq K_N^\mu$
(the unit ball of the smaller form is larger,
so the Christoffel kernel is larger).

\begin{prop}[kernel Loewner]\label{prop:kernel}
If the tilde chain is positive through $N$,
then for every finite atomic signed border
$\sigma=\sum_a w_a\delta_{t_a}$,
\[
  \|b\|^2
  \;=\;
  w^\top K_N^\mu w
  \;\le\;
  w^\top K_N^{\mu-\nu} w
  \;=\;
  S.
\]
Over $\mathbb{Q}$: $\mu=(2,1)$ at $(0,1)$,
$\nu=(1)$ at $1$, constants,
$K^\mu=\tfrac13\le K^{\mu-\nu}=\tfrac12$,
Dirac at $0$ has ratio $\tfrac23$.
At $k_z=9$, the Rayleigh is $0.7264$;
the operator norm of the change-of-basis on
$\mathrm{col}(\Psi)$ is $1.0000$.
\end{prop}

The naive direction
``$L^2(\mu)\ge L^2(\mu-\nu)$ $\Rightarrow$
$\mu$-Parseval $\ge$ tilde-Parseval''
is the \emph{wrong} way for measures:
a larger form has a \emph{smaller} kernel.
That is why $\|b\|^2\le S$, not the reverse.

The full Parseval sum over all $\mu$-degrees
is not a tame $\sigma$-norm:
$\sigma$ sits on border atoms, typically off
the $\mu$-support, so $K_M^\mu(t,t)$ explodes
as $M$ approaches the discrete $\mu$-degree.

\section{Weighted Cauchy--Schwarz / Gram}\label{sec:gram}

Termwise $w_k=b_k^2/\rho_k$ remains unbounded
(r424).
The exact cross-identity is not
$b_k=c\,F_k$ per index; it is the
change-of-basis $b=Gf$ on $P_{<N}$,
with $f_k=F_k/\sqrt{\eta_k}$ and $\|f\|^2=S$.
$\|G\|_{\mathrm{op}}=1$ on $\mathrm{col}(\Psi)$
--- a construction object (the identity map
between the two norms).
The true window border is \emph{not} the
maximizer: random borders have Rayleigh
$\sim 0.02$; the one-atom maximizer-near
direction at $k_z=9$ has $0.75$;
$k=5$ is the tight living ray $0.799$.
A uniform $\|G\|_{\mathrm{op}}\le 0.9$ is
false ($\|G\|_{\mathrm{op}}=1$).
The $0.80$ census is the Rayleigh of
\emph{this} $\sigma$, not an operator bound.

\section{Sharp form, reach, Lean}\label{sec:reach}

The sharp theorem is $\|b\|^2\le S$.
On any window with $q_N<1$,
$S<B_w$ and therefore $\gamma<1$.
The $42$-rung $q_N<1$ census
($q_N\in[0.0015,0.9805]$) plus
Proposition~\ref{prop:kernel} gives
$\gamma<1$ and, with the r423 formula,
$\mathrm{den}<2$ as a composed finite
statement on those rungs.
$\Sigma$-boundedness and the reserve floor
stay census/evidence; the cofinal edge
is still conditional on cofinal $q_N<1$.

Lean sketch (not this round):
a named \texttt{Prop} on a finite-dimensional
real inner-product space,
\texttt{inner\_mu $\ge$ inner\_mutilde}
(certificate: two Gram matrices,
Cholesky of the difference
$=\sum_j v_j p(y_j)^2$),
then the kernel Loewner and the pairing
inequality.  $\mathbb{Q}$-toy as an instance.

\section{Kills}\label{sec:kills}

\begin{center}
\begin{tabular}{llll}
\toprule
world / mutant & $\|b\|^2\le S$ & $q_N<1$ & $\gamma<1$ \\
\midrule
core-$42$ + EXT + selected & yes & yes & SATZ (composed) \\
$k=5$ (tight ray $0.799$) & yes & $0.023$ & kernel slack \\
dead $\chi$ & yes ($0.62$) & \textbf{no} ($1.21$, $1.33$) & slack, not $q_N$ \\
random border & yes ($\sim 0.02$) & (chain ok) & far from max \\
atom reorder & yes & yes & measure, not order \\
scramble & n/a & n/a & tilde form not definite \\
unnorm $b$ & n/a & --- & $\gamma=2.51>1$ \\
\bottomrule
\end{tabular}
\end{center}

Dead $\chi$ is the named witness that
$\|b\|^2\le S$ does \emph{not} by itself
give $\|b\|^2\le B_w$ when $q_N>1$.

\section{What is a theorem}\label{sec:honest}

Proposition~\ref{prop:kernel} is a theorem.
Lemma~\ref{lem:cross} is not cofinal:
it is this SATZ plus $q_N<1$.
On the named $42$-rung census the composition
is a theorem and $\mathrm{den}<2$ follows.
No RH claim.

\end{document}
